\documentclass[10pt]{article}

\usepackage[utf8]{inputenc}
\usepackage[margin=0.85in]{geometry}
\usepackage{booktabs}
\usepackage{amsmath}
\usepackage{microtype}
\usepackage{xcolor}
\usepackage{hyperref}
\usepackage{listings}

\hypersetup{
  pdftitle={flower: Growing Agentic Computational Work into Reproducible Protocols},
  pdfauthor={Zhanghao Zhouyin},
  colorlinks=true,
  linkcolor=blue!55!black,
  citecolor=blue!55!black,
  urlcolor=blue!55!black
}

\definecolor{codebg}{RGB}{247,247,247}
\lstdefinelanguage{yaml}{
  sensitive=false,
  comment=[l]{\#},
  morestring=[b]",
  morestring=[b]'
}
\lstset{
  basicstyle=\ttfamily\small,
  backgroundcolor=\color{codebg},
  frame=single,
  rulecolor=\color{black!15},
  breaklines=true,
  columns=fullflexible,
  showstringspaces=false,
  aboveskip=0.7em,
  belowskip=0.7em
}

\newcommand{\flower}{\textsc{flower}}

\title{\flower{}: Growing Agentic Computational Work into Reproducible Protocols}
\author{Zhanghao Zhouyin\\McGill University}
\date{}

\begin{document}

\maketitle

\begin{abstract}
Coding agents can now carry out substantial computational studies, but the record they leave is often
harder to assess than the final result. Commands are tried in terminals, scripts are revised in place,
long jobs outlive conversations, and the path from an initial question to a reported number is lost among
abandoned attempts. Conventional workflow systems provide reliable execution once a pipeline has been
specified, but exploratory work rarely begins with a finished pipeline. We present \flower{}, a lightweight,
file-first workflow tool designed around this gap. A study begins immediately as an empty, inspectable run.
The agent adds and revises steps while an append-only journal records attempts, results, decisions, and plan
changes. Successful steps are stabilized as scripts with declared inputs, outputs, files, environments, and
compute resources. At the end, \flower{} extracts the dependency closure behind a selected result as a
re-runnable protocol with expected outputs. The same representation supports local processes, remote
workstations, Slurm clusters, human gates, crash recovery, and a visual flow map. In this way, the exploratory
history remains auditable while the final deliverable becomes smaller, deterministic, and reusable.
\end{abstract}

\section{The missing object between a conversation and a pipeline}

An agent asked to reproduce a scientific result does not normally know the final workflow in advance. It
first finds data, tests an interpretation, writes a small program, discovers that a dependency is missing,
changes the program, launches a larger calculation, and revises the analysis when the result exposes a
mistake. This is not an exception to scientific work; it is how the method is found.

The difficulty is that the natural record of this process is a conversation plus a working directory. The
conversation explains some decisions but is too long to audit as an execution trace. The directory contains
the latest scripts but not necessarily the versions that produced the reported result. A scheduler knows that
a job completed but not why the job was launched or which later conclusion depended on it. When an agent
returns a number, a human reviewer must reconstruct these relationships by hand.

Mature workflow systems such as Snakemake~\cite{koster2012snakemake} and
Nextflow~\cite{ditommaso2017nextflow} make declared pipelines portable and repeatable. Provenance-oriented
systems such as AiiDA~\cite{huber2020aiida} preserve rich graphs of calculations and data. These systems solve
important parts of the reproducibility problem. The design point of \flower{} is earlier in the lifecycle: the
period in which the workflow does not yet exist. Its central question is therefore not only ``How should this
pipeline run?'' but also ``How can the process of discovering the pipeline become the evidence for the final
one?''

\section{A study that grows}

\flower{} treats a study as three related objects. A \emph{plan} is the current human-readable description of
the work. A \emph{run} is the plan together with its immutable history. A \emph{protocol} is the result-bearing
part of a successful run, prepared for another person or machine to execute again.

The run is created before exploration begins. It may start with no nodes at all. Each computation that the
agent relies on---including a download, a quick numerical check, or a failed pilot---is added as a step. A
step states what it is meant to establish, what command it executes, which earlier results it uses, and which
outputs or files it promises to produce. References to an earlier output also create dependency edges. A
small plan can therefore remain close to the commands a researcher would have typed directly:

\begin{lstlisting}[language=yaml,caption={A minimal \flower{} plan.}]
flower: 1
id: silicon-bands
title: Silicon band structure
nodes:
  - id: relax
    kind: shell
    description: Relaxes the cell; energy is the convergence result.
    run: python3 relax.py > "$FLOWER_OUTPUTS"
    outputs: {energy: number, converged: boolean}

  - id: bands
    kind: shell
    description: Computes and plots bands at the relaxed geometry.
    when: "${relax.outputs.converged}"
    run: python3 bands.py > "$FLOWER_OUTPUTS"
    files: {plot: bands.png}
\end{lstlisting}

This simplicity is deliberate. The agent remains outside the workflow engine and uses its existing coding
environment. \flower{} does not provide an LLM loop, prompt language, or tool registry. The agent decides what
to try; \flower{} executes the durable effect. Today the executable vocabulary is primarily a shell step and a
human gate. Local programs, Python modules, simulation packages, and remote jobs all fit behind the same shell
contract. This keeps the recorded workflow independent of a particular agent vendor or model.

Exploration does not require history to be rewritten. New steps and changes to unfinished steps become new
plan generations. A change to completed work is explicit: the old attempt remains in the journal, the changed
node is run again, and affected descendants become stale. Side studies and previews can remain in the run
without contaminating the final protocol. They are evidence of how the method was found, not necessarily part
of how the result must be reproduced.

\begin{figure}[t]
  \centering
  \setlength{\tabcolsep}{5pt}
  \renewcommand{\arraystretch}{1.45}
  \begin{tabular}{c@{\quad$\longrightarrow$\quad}c@{\quad$\longrightarrow$\quad}c}
    \fbox{\begin{minipage}{0.23\linewidth}\centering
      \textbf{Explore}\\
      add probes, fail, revise
    \end{minipage}}
    &
    \fbox{\begin{minipage}{0.23\linewidth}\centering
      \textbf{Stabilize}\\
      scripts, contracts, environments
    \end{minipage}}
    &
    \fbox{\begin{minipage}{0.23\linewidth}\centering
      \textbf{Deliver}\\
      export, rerun, compare
    \end{minipage}}
  \end{tabular}
  \caption{The run retains exploratory history, while protocol export keeps only the steps behind a selected
  result.}
  \label{fig:lifecycle}
\end{figure}

\section{Durability as a file format}

The source of truth for a run is an append-only JSON Lines journal. Every event carries a contiguous sequence
number, timestamp, actor label, and closed event type. Appends are serialized with a file lock and flushed to
disk. Idempotency keys make retried writes harmless. A torn final line can be discarded after a crash, whereas
corruption in the middle of the history is reported rather than guessed around. The visible run state is a
pure fold over this journal; derived status files and views may be deleted and rebuilt.

Execution advances in short, re-entrant passes. A pass folds the journal, reconciles running work, records
completed attempts, and starts newly ready nodes. Local commands run below detached supervisors. Remote work
is either a detached process on a workstation or a batch job submitted to Slurm. Thus an agent session, UI,
or foreground command does not need to remain alive while a calculation runs. A later status command can
continue the same run.

Remote submission is treated as an irreversible boundary. Before submission, \flower{} records an intent with
a deterministic key. Submission, remote markers, scheduler lookup, and an in-job ownership guard are arranged
so that recovery favors reattachment over duplicate execution. Polling distinguishes scheduler delay,
transport failure, preemption, node failure, memory exhaustion, timeout, cancellation, and a genuinely lost
job. These distinctions matter because blindly repeating a scientific calculation is not always safe or
useful.

Each successful attempt records two identities. The declaration hash describes what the node does, including
the content of referenced local source files; the input hash describes resolved inputs and upstream output
and file hashes. A result can be reused only when both identities match:
\[
  \operatorname{reuse}(n) \iff
  h_{\mathrm{decl}}^{\mathrm{new}}(n)=h_{\mathrm{decl}}^{\mathrm{old}}(n)
  \;\land\;
  h_{\mathrm{input}}^{\mathrm{new}}(n)=h_{\mathrm{input}}^{\mathrm{old}}(n).
\]
Scheduler resources and descriptive text are excluded from the declaration identity, so increasing memory or
improving an explanation does not invalidate a scientific result. A deliberate rerun can still force the
selected node while allowing unchanged descendants to be reused.

Software environments follow the same transition from exploration to stabilization. Commands used to explore
a target machine can be collected into a small recipe consisting of setup, activation, and verification
scripts. Freezing the recipe assigns it a content hash and snapshots its files. A fresh replay demonstrates
that the recipe installs into an empty prefix. Changing the recipe produces a new prefix and changes the cache
identity of every dependent step.

\section{A flow map for human oversight}

The journal is optimized for correctness, not for reading. \flower{} therefore projects it into a visual map of
the current directed acyclic graph. Nodes show their purpose before their machinery: the description, result,
declared files, attempts, logs, and failure details are available from the same view. A timeline presents the
events in order, and a plan-history view shows each generation and amendment. Gates place a human decision in
the graph itself rather than leaving approval in a separate conversation.

This human view is important because observability and reproducibility are different properties. A workflow
may be technically replayable and still be difficult to judge. Requiring each step to say what it establishes
gives the graph an explanatory layer. Conversely, the description alone is not trusted as execution history:
the journal, output contracts, and file hashes record what actually happened. The model resembles the
distinction in provenance standards between activities, entities, and responsible agents~\cite{w3cprov}, but
uses a compact, project-local representation during active work.

Human authority is represented through plan approvals, amendment decisions, and gate nodes. The current
implementation records the declared actor associated with each action; it does not authenticate that identity.
This is sufficient for collaboration on a trusted machine but not yet for institutional audit or regulatory
use.

\section{From a run to a protocol}

When a result is ready, the user selects one or more successful nodes. \flower{} computes their ancestor closure
and writes three files. \texttt{protocol.yaml} contains the selected nodes and everything they depend on;
\texttt{expected.json} contains the structured results observed in the original run; and
\texttt{PROTOCOL.md} explains the required inputs, machines, software recipes, steps, and reproduction command.
Exploratory branches that do not contribute to the selected result disappear by construction.

The protocol is executed through the same engine as the exploratory run. Afterwards, a comparison command
checks the new structured outputs against the recorded expectations, using numerical tolerances where
appropriate. Protocol extraction is therefore not a separate packaging convention placed on top of the
workflow. It is the final operation of the workflow model: historical evidence is reduced to a result-bearing
recipe without erasing the original history.

\section{Experience from computational science}

The current implementation has been exercised through sixteen tracked reproductions spanning condensed
matter physics, electronic structure, quantum chemistry, and molecular simulation. The workloads include
exact diagonalization, quantum Monte Carlo, density-functional calculations, density-matrix renormalization,
molecular dynamics, and coupled-cluster calculations. They run across the local machine, remote workstations,
and Slurm-style execution, with fan-out campaigns ranging from small checks to tens of independent jobs.

These studies were used as design probes rather than as a single performance benchmark. They exposed failure
modes that synthetic happy-path examples rarely reveal: a foreach campaign whose item list changes after an
upstream rerun; a process that survives its driver; a scheduler submission whose reply is lost; a completed
remote job whose result download temporarily fails; a checkpoint that should survive an automatic retry but
not an intentional rerun; and a descriptive edit that should not repeat days of computation. The repository
currently contains 359 automated test functions, including fake-scheduler fault injection and crash-window
tests. More importantly, observed failures are retained as a bug diary and converted into regression tests.

The experiments also support the main usability claim. A study can begin before its final graph is known, grow
while computations are running, and later yield a smaller protocol that has been rerun from a fresh checkout.
The value of the tool is therefore not measured only by successful scheduling. It is measured by whether a
person who was absent during the exploration can understand which steps support the final claim and can ask a
new machine to perform them again.

\section{Limitations and next steps}

\flower{} is currently a project-local tool. Runs in separate clones are independent, actors are declared rather
than authenticated, and there is no shared service for merging histories. The journal is crash-safe but not
yet tamper-evident. Signed hash chains and authenticated approvals are natural extensions for results that must
cross organizational boundaries.

Reproducibility also remains conditional on the quality of the recipe. \flower{} can guarantee that a frozen
recipe is content-addressed and replayed consistently, but it cannot make an unpinned package installer
deterministic. The present protocol comparison emphasizes structured JSON outputs; richer deliverables require
per-output validation policies and explicit checks for large files, plots, and stochastic results. A future
export manifest should bind the protocol to the exact source tree, environment recipes, declared artifacts,
and git state that produced the selected results.

Finally, the current representation asks exploratory computations to become steps early. This is useful
discipline, but lightweight probe events and an explicit ``promote to stable step'' operation could better
represent the boundary between an observation used during exploration and a node intended for the final
protocol.

\section{Conclusion}

Agentic computational work creates a new provenance problem: execution becomes faster while human
reconstruction becomes harder. \flower{} addresses this problem by making the evolving workflow, rather than
the conversation, the durable object. It allows a study to grow before its structure is known, preserves the
failed and revised path by which that structure was found, and extracts the successful dependency closure as a
re-runnable protocol. The result is a small bridge between open-ended agency and deterministic delivery: the
agent is free to explore, the machine can recover, and the human receives both a readable flow map and a
recipe that can be checked.

\paragraph{Availability.}
\flower{} is open-source software under the Apache License 2.0. The implementation and benchmark workflows are
available at \url{https://github.com/floatingCatty/flower}.

\begin{thebibliography}{9}

\bibitem{koster2012snakemake}
J. Koster and S. Rahmann,
``Snakemake---a scalable bioinformatics workflow engine,''
\emph{Bioinformatics}, vol. 28, no. 19, pp. 2520--2522, 2012.
\href{https://doi.org/10.1093/bioinformatics/bts480}{doi:10.1093/bioinformatics/bts480}.

\bibitem{ditommaso2017nextflow}
P. Di Tommaso, M. Chatzou, E. W. Floden, P. P. Barja, E. Palumbo, and C. Notredame,
``Nextflow enables reproducible computational workflows,''
\emph{Nature Biotechnology}, vol. 35, pp. 316--319, 2017.
\href{https://doi.org/10.1038/nbt.3820}{doi:10.1038/nbt.3820}.

\bibitem{huber2020aiida}
S. P. Huber et al.,
``AiiDA 1.0, a scalable computational infrastructure for automated reproducible workflows and data provenance,''
\emph{Scientific Data}, vol. 7, article 300, 2020.
\href{https://doi.org/10.1038/s41597-020-00638-4}{doi:10.1038/s41597-020-00638-4}.

\bibitem{w3cprov}
L. Moreau and P. Missier, eds.,
``PROV-DM: The PROV Data Model,'' W3C Recommendation, 30 April 2013.
\url{https://www.w3.org/TR/prov-dm/}.

\end{thebibliography}

\end{document}
